\uuid{2kfm}
\exo7id{3604}
\titre{exo7 3604}
\auteur{quercia}
\organisation{exo7}
\datecreate{2010-03-10}
\isIndication{false}
\isCorrection{false}
\chapitre{Réduction d'endomorphisme, polynôme annulateur}
\sousChapitre{Autre}
\cpge{mpsi-mp2i.espaces-vectoriels.a, pcsi.espaces-vectoriels.a, ptsi.geometrie-espace.c, mp-mpi.reduction-endomorphismes.d, pc.reduction.c, psi.reduction-endomorphismes.c, pt.reduction-endomorphismes.c, mp-mpi.reduction-endomorphismes.b, pc.reduction.a, psi.reduction-endomorphismes.a, pt.reduction-endomorphismes.a, bcpst2.valeurs-propres-vecteurs-propres.b}
\module{Algèbre}
\niveau{L2}
\difficulte{}

\contenu{
\question{
Soit $E$ un $ K$-ev de dimension finie et $u \in \mathcal{L}(E)$ diagonalisable.
Pour $\lambda \in \text{spec}(u)$, on note $E_\lambda = \mathrm{Ker}(u-\lambda\mathrm{id})$
et $F_\lambda = \Im(u-\lambda\mathrm{id})$.
Montrer que $E_\lambda \oplus F_\lambda = E$.
}
}
